\section{Length-Weighting Derivation}
\label{app:normalization}

Fix a rollout batch, its valid-token masks, and the domain weights
$\lambda_d$. Any PG advantages are detached. Use the response mean
$\ell_i$ and response gradient $g_i=\nabla_\theta\ell_i$ from
Section~\ref{sec:normalization}. For $n_d=|\mathcal B_d|$, define mean
response length $\bar T_d=n_d^{-1}\sum_{i\in\mathcal B_d}|y_i|$ and mean
response gradient $\bar g_d=n_d^{-1}\sum_{i\in\mathcal B_d}g_i$.
The empirical length--gradient covariance uses the same $1/n_d$
normalization:
\begin{equation}
\operatorname{Cov}_d(|y|,g)
=\frac{1}{n_d}\sum_{i\in\mathcal B_d}
(|y_i|-\bar T_d)(g_i-\bar g_d),\qquad
\frac{\sum_{i\in\mathcal B_d}|y_i|g_i}
     {\sum_{i\in\mathcal B_d}|y_i|}
=\bar g_d+\frac{\operatorname{Cov}_d(|y|,g)}{\bar T_d}.
\label{eq:norm_covariance}
\end{equation}
To verify the identity, expand the covariance as
$n_d^{-1}\sum_i|y_i|g_i-\bar T_d\bar g_d$ and divide by $\bar T_d$.
Global token averaging combines the right-hand side using each domain's
observed token share; domain token averaging uses $\lambda_d$ instead.
Domain response averaging gives $\sum_d\lambda_d\bar g_d$.
Thus the change in the domain coefficient and the within-domain
length--gradient covariance are distinct, even with identical rollouts.
Multiplying a domain's globally averaged token terms by its target weight
divided by its observed token share gives the domain-token rule exactly.

These are finite-batch identities. The expectation of a sample ratio
is generally not the ratio of population expectations; replacing it
with $\mathbb E[|y|g]/\mathbb E[|y|]$ requires an appropriate large-batch
limit. The identities also do not differentiate through the response
sampling distribution, masks, or response lengths.

\section{Measurement Details and Optional Follow-up}

\subsection{Teacher-conditioned masks at a common joint checkpoint}
\label{app:records}

The diagnostic record contains the joint student checkpoint, optimizer
state, teacher identity, routed domain, prompt IDs, sampled response
and mask, loss reduction, and gradient accumulation denominator.
Each teacher-conditioned gradient and proposed step is computed from
the same restored student state. They describe the parameters most affected by a teacher at that
checkpoint, not a unique historical attribution of the cumulative joint delta.

Record raw and clipped gradients separately. Retain the actual optimizer
step and its momentum and decay settings. A zero-gradient optimizer
branch can identify motion caused by stored momentum or weight decay
when those components dominate an overlap comparison. Prefer FP32
master-weight differences for cumulative updates when available;
converting already rounded BF16 checkpoints to FP32 does not recover
unrecorded updates. Record exclusions, thresholds, relative scales,
tie rules, and per-layer selected-parameter counts.

Comparisons selecting an equal fraction of parameters accompany
absolute-threshold selections.
A stratified random baseline preserves layer and source-weight-magnitude
bins when feasible. Degenerate or tie-dominated gradients are marked;
ranking a zero vector does not define a meaningful learned subnetwork.
If both selections are empty, report overlap as undefined and exclude
that comparison from averages. If only one is empty, overlap is zero.

\subsection{Teacher distribution difference}
\label{app:teacher_distance}

For teachers $d$ and $e$ at the same prefix, the Jensen--Shannon divergence is
\begin{equation}
\operatorname{JS}(\pi_{\phi_d},\pi_{\phi_e})
=\tfrac12 D_{\mathrm{KL}}(\pi_{\phi_d}\|m)
+\tfrac12 D_{\mathrm{KL}}(\pi_{\phi_e}\|m),
\qquad m=\tfrac12(\pi_{\phi_d}+\pi_{\phi_e}).
\label{eq:teacher_js}
\end{equation}
It lies between zero and $\log 2$ with natural logarithms and is zero
when the two next-token distributions agree. Average this quantity
across the same prefix bank and domain weights for every pair.
Replacing one teacher distribution with $\pi_\theta$ gives the
teacher--student comparison. Truncated estimates must report their
retained probability mass and must not be presented as exact
full-vocabulary divergences.

\subsection{Optional analysis if the loss comparison shows a stable difference}
\label{app:optional_loss}

If a stable difference warrants further analysis, resample next actions
on a fixed student-prefix bank and compare sample counts, an
action-independent detached baseline, and the exact full-vocabulary
gradient. These controls test whether the effect tracks sampling variance. A retained-head plus sampled-tail estimator can distinguish
truncation from the full reverse-KL reference
\citep{oh2026vopd,zhang2026emapg,yu2026denseupdates}.

The independent action draws do not extend the response; prefix selection
and weights remain fixed. Original rollout tokens with eventual
response-length weights need not satisfy the same sampling identity. Forward KL, teacher versus student top-$k$, or within-set
renormalization are objective changes and are labeled separately.


\subsection{An optional amplitude calibration for teacher distance}
\label{app:teacher_controls}

A simple artificial target can help distinguish gradient magnitude from
parameter selection. At a fixed prefix $(x,y_{<t})$, write
$\pi_0=\pi_{\theta_0}$ for the reference student and freeze the teacher
$\pi_{\phi_d}$. Define the normalized interpolated distribution
\begin{equation}
\pi_\alpha(v)=
\frac{\pi_0(v)^{1-\alpha}\pi_{\phi_d}(v)^\alpha}
     {\sum_{v'}\pi_0(v')^{1-\alpha}\pi_{\phi_d}(v')^\alpha},
\qquad
\left.\nabla_\theta D_{\mathrm{KL}}(\pi_\theta\|\pi_\alpha)
\right|_{\theta_0}
=\alpha\left.\nabla_\theta D_{\mathrm{KL}}
(\pi_\theta\|\pi_{\phi_d})\right|_{\theta_0}.
\label{eq:tilt_gradient}
\end{equation}
All probabilities condition on the same prefix. The log-ratio is an
$\alpha$-scaled original log-ratio plus a constant across vocabulary tokens;
the constant vanishes because $\sum_v\nabla_\theta\pi_\theta(v)=0$.
For $\alpha>0$, positive scaling preserves the ranking of raw-gradient
magnitudes while possibly changing absolute-threshold sparsity.
This artificial calibration does not replace real-teacher comparisons
or imply matched downstream capability.

\subsection{Reproducibility records}

Record exact student, teacher, tokenizer, data, and evaluator revisions;
teacher training lineages; prompt schedules and seeds; decoding,
EOS handling, cap, invalid responses, and loss masks; learning-rate
schedule, optimizer moments, precision, clipping, and decay; candidate
sets, probability normalization, teacher/student retained mass, and
advantage handling. Keep per-prompt outputs, source measurements for
every plot, and actual time and memory costs.

Distinguish training seeds, diagnostic batch replicates, and evaluation
resamples. Pairs sharing a teacher are dependent, and prompt-level
uncertainty does not establish population-level teacher relationships.
Historical synthetic GPAS calculations are not empirical evidence for these studies.
